\documentclass{article}

\usepackage{amsmath,amssymb}
\usepackage{xurl}
\usepackage[hidelinks]{hyperref}

% Shared commands for notes in docs/paper-gaps/.

\DeclareUnicodeCharacter{2080}{\ifmmode{}_{0}\else\textsubscript{0}\fi}
\DeclareUnicodeCharacter{2081}{\ifmmode{}_{1}\else\textsubscript{1}\fi}
\DeclareUnicodeCharacter{2082}{\ifmmode{}_{2}\else\textsubscript{2}\fi}
\DeclareUnicodeCharacter{2099}{\ifmmode{}_{n}\else\textsubscript{n}\fi}

\newcommand{\Lean}{\textsc{Lean}}
\ExplSyntaxOn
\NewDocumentCommand{\leanid}{m}{
  \ifmmode
    \textsf{\detokenize{#1}}
  \else
    \group_begin:
    \urlstyle{sf}
    \tl_set:Nn \l_tmpa_tl {#1}
    \tl_replace_all:Nnn \l_tmpa_tl {\_} {_}
    \exp_args:NV \path \l_tmpa_tl
    \group_end:
  \fi
}
\ExplSyntaxOff
\providecommand{\tr}{\operatorname{tr}}
\newcommand{\tnleanIssueBase}{https://github.com/LionSR/TNLean/issues/}
\newcommand{\tnleanPrBase}{https://github.com/LionSR/TNLean/pull/}
\newcommand{\tnleanDocBase}{https://sirui-lu.com/TNLean/docs/}
\newcommand{\ghissue}[1]{\href{\tnleanIssueBase#1}{GitHub issue \##1}}
\newcommand{\ghpr}[1]{\href{\tnleanPrBase#1}{PR \##1}}
\newcommand{\doclink}[1]{\href{\tnleanDocBase#1.html}{\path{#1}}}

% Dirac notation and the identity operator, as in blueprint/src/macros/common.tex.
% \providecommand keeps note-local definitions authoritative.
\usepackage{dsfont}
\providecommand{\ket}[1]{|#1\rangle}
\providecommand{\bra}[1]{\langle #1|}
\providecommand{\braket}[2]{\langle #1 | #2 \rangle}
\providecommand{\ketbra}[2]{|#1\rangle\!\langle #2|}
\providecommand{\Id}{\mathds{1}}
\providecommand{\one}{\mathds{1}}
\providecommand{\C}{\mathbb{C}}
\providecommand{\MN}[1]{M_{#1}(\C)}
\providecommand{\MD}{M_D(\C)}

% External referencing for paper sources.  The build helper writes sanitized
% auxiliary files under build/paper-aux/ and excludes duplicated source labels.
\usepackage{xr}

% Import a generated paper auxiliary file with a caller-supplied label prefix.
% The candidate paths support builds from docs/paper-gaps/, the repository root,
% and build/paper-aux/ itself.
\newcommand{\importpaperaux}[2]{%
  \IfFileExists{../../build/paper-aux/#2.aux}{%
    \externaldocument[#1]{../../build/paper-aux/#2}%
  }{%
    \IfFileExists{build/paper-aux/#2.aux}{%
      \externaldocument[#1]{build/paper-aux/#2}%
    }{%
      \IfFileExists{#2.aux}{%
        \externaldocument[#1]{#2}%
      }{}%
    }%
  }%
}

% General external reference macros
\newcommand{\paperref}[2]{\ref{#1-#2}}
\newcommand{\papereqref}[2]{(\ref{#1-#2})}

% CPSV16 source (1606.00608 / MPDO-22-12-17-2.tex) external document and shortcuts
\importpaperaux{cpsv16-}{1606.00608/MPDO-22-12-17-2-tnlean}
\newcommand{\cpsvref}[1]{\paperref{cpsv16}{#1}}
\newcommand{\cpsveqref}[1]{\papereqref{cpsv16}{#1}}

% Verdict marker: \gapnote{<kind>}{<status>}. Kind is the policy
% classification (clarification, local-correction, scope-restriction,
% unfaithful, false-source, open-gap); status is open, wip (elimination
% underway), resolved, or historical. Machine-read by the site index;
% prints a verdict line.
\newcommand{\gapnote}[2]{\par\noindent\textbf{Verdict:} #1 (#2).\par}


\title{Fixed-Length Chain Equality Does Not Imply Combined-Tensor MPV Equality}
\author{}
\date{2026-07-20}

\begin{document}
\maketitle
\gapnote{unfaithful}{resolved}

\section{The incompatible statements}

Let $\mathbf A=(A_0,\ldots,A_{n-1})$ be a site-dependent closed chain. Equality
of two states at the fixed chain length $n$ tests only cyclic products that
contain one matrix from each site in cyclic order. By contrast, MPV equality
for the tensor formed by combining the site and physical indices tests every
word over the enlarged alphabet. Such a word may repeat one site while
omitting all other sites. Fixed-length state equality therefore does not imply
MPV equality for the combined tensors.

A scalar example already proves the incompatibility. Take physical dimension
one, bond dimension one, and chain length three, and define
\begin{align}
    A_0
        &=
        A_1
        =
        A_2
        =
        1,
    \label{eq:tnlean_algebraic_ft_same_state_combined_mpv_gap_first_chain}\\
    B_0
        &=
        2,
    \qquad
    B_1
        =
        \frac{1}{2},
    \qquad
    B_2
        =
        1.
    \label{eq:tnlean_algebraic_ft_same_state_combined_mpv_gap_second_chain}
\end{align}
Every local tensor in both chains is injective. Their closed-chain
coefficients agree because
\begin{align}
    A_0A_1A_2
        &=
        1,
    \label{eq:tnlean_algebraic_ft_same_state_combined_mpv_gap_first_chain_product}\\
    B_0B_1B_2
        &=
        2\cdot\frac{1}{2}\cdot 1
        =
        1.
    \label{eq:tnlean_algebraic_ft_same_state_combined_mpv_gap_second_chain_product}
\end{align}
Nevertheless, the length-one word that selects site zero has different
coefficients in the combined tensors:
\begin{align}
    A_0
        &=
        1
        \neq
        2
        =
        B_0.
    \label{eq:tnlean_algebraic_ft_same_state_combined_mpv_gap_length_one_mismatch}
\end{align}
Thus the equality established by
(\ref{eq:tnlean_algebraic_ft_same_state_combined_mpv_gap_first_chain_product})
and
(\ref{eq:tnlean_algebraic_ft_same_state_combined_mpv_gap_second_chain_product})
does not imply equality of the combined-tensor MPV families.

Consequently, the former abstract same-state reduction assertion has no valid
construction, even when injectivity of the second chain is added.%
\footnote{Former formal declaration:
    \path{MPSChainTensor.SameStateBridgeHyp}, previously in
    \path{TNLean/MPS/Chain/SameStateBridge.lean}.}
In particular, this assertion is not the blocking argument in
Theorem~\textup{thm:inj\_MPS}, lines~688--725, of
Ref.~\cite{Molnar2018NormalPEPS}.

\section{The source-faithful route}

The source theorem assumes that every local tensor in both chains is
injective. From equality of the single closed-chain state, it derives one
invertible gauge for each bond
\cite[Theorem~\textup{thm:inj\_MPS}, lines~688--725]{Molnar2018NormalPEPS}.
Chapter~24 of the blueprint formalizes this statement directly.%
\footnote{Formal declaration:
    \leanid{TNLean.PEPS.fundamentalTheorem_injectiveMPSChain_of_sameState}
    in \doclink{TNLean/PEPS/CycleMPSChainOverlapCapstone}.}
The proof does not introduce all-length MPV equality for a tensor with
combined site and physical indices.

A related phase obstruction occurs for constant chains. Equality at one length
determines the tensor only up to a similarity and a scalar whose $n$-th power
is one. To see this in bond dimension one, choose a nontrivial $n$-th root of
unity $\omega$ and set
\begin{align}
    A
        &=
        1,
    \qquad
    B
        =
        \omega,
    \qquad
    \omega^n
        =
        1,
    \qquad
    \omega
        \neq
        1.
    \label{eq:tnlean_algebraic_ft_same_state_combined_mpv_gap_root_of_unity_example}
\end{align}
The length-$n$ coefficients agree:
\begin{align}
    A^n
        &=
        1
        =
        \omega^n
        =
        B^n.
    \label{eq:tnlean_algebraic_ft_same_state_combined_mpv_gap_fixed_length_phase_equality}
\end{align}
Scalar conjugation cannot transform $A=1$ into $B=\omega$, so a phase-free
conclusion from one fixed length is false. The all-length hypothesis in the
translation-invariant reduction excludes this phase and defines a separate,
strictly stronger theorem.%
\footnote{Formal declaration:
    \leanid{MPSChainTensor.ti_reduction_corollary} in
    \doclink{TNLean/MPS/Chain/TranslationInvariance}.}

\section{Verdict}

The classification is an unfaithful theorem. The abstract same-state reduction
hypothesis asserted a false implication that is not supplied by
Ref.~\cite{Molnar2018NormalPEPS}; the hypothesis and its consumers have
therefore been removed.

The fixed-length Fundamental Theorem for injective chains is represented only
by the direct, source-faithful theorem in Chapter~24. The Chapter~23 results
based on combined-tensor MPV equality remain explicitly all-length results.
They are not consequences of equality of one fixed-length state.

\bibliographystyle{alpha}
\bibliography{references}

\end{document}
